% 前言
\chapter{Preface}

% 增大字体并调整行距
\large
\setlength{\parskip}{0.3\baselineskip} % 增加段落间距
\setlength{\parindent}{2em} % 添加首行缩进
\vspace{6.5cm}
This book is intended as a text for a one- or two-semester introduction to topology, 
at the senior or first-year graduate level.

The subject of topology is of interest in its own right, and it also serves to lay 
the foundations for future study in analysis, in geometry, and in algebraic topology. 
There is no universal agreement among mathematicians as to what a first course in 
topology should include; there are many topics that are appropriate to such a course, 
and not all are equally relevant to these differing purposes. In the choice of material 
to be treated, I have tried to strike a balance among the various points of view.

\noindent\textbf{Prerequisites. }There are no formal subject matter prerequisites 
for studying most of this book. I do not even assume the reader knows much set theory. 
Having said that, I must hasten to add that unless the reader has studied a bit of 
analysis or ``rigorous calculus'', much of the motivation for the concepts introduced 
in the first part of the book will be missing. Things will go more smoothly if he or 
she already has had some experience with continuous functions, open and closed sets, 
metric spaces, and the like, although none of these is actually assumed. In Part II, 
we do assume familiarity with the elements of group theory.

Most students in a topology course have, in my experience, some knowledge of the 
foundations of mathematics. But the amount varies a great deal from one student to 
another. Therefore, I begin with a fairly thorough chapter on set theory and logic. 
It starts at an elementary level and works up to a level that might be described as 
``semi-sophisticated.'' It treats those topics (and only those) that will be needed 
later in the book. Most students will already be familiar with the material of the 
first few sections, but many of them will find their \textit{expertise} disappearing 
somewhere about the middle of the chapter. How much time and effort the instructor 
will need to spend on this chapter will thus depend largely on the mathematical 
sophistication and experience of the students. Ability to do the exercises fairly 
readily (and correctly!) should serve as a reasonable criterion for determining 
whether the student's mastery of set theory is sufficient for the student to begin 
the study of topology.

Many students (and instructors!) would prefer to skip the foundational material of 
Chapter 1 and jump right in to the study of topology. One ignores the foundations, 
however, only at the risk of later confusion and error. What one \textit{can} do is 
to treat initially only those sections that are needed at once, postponing the 
remainder until they are needed. The first seven sections (through countability) 
are needed throughout the book; I usually assign some of them as reading and lecture 
on the rest. Sections 9 and 10, on the axiom of choice and well-ordering, are not 
needed until the discussion of compactness in Chapter 3. Section 11, on the maximum 
principle, can be postponed even longer; it is needed only for the Tychonoff theorem 
(Chapter 5) and the theorem on the fundamental group of a linear graph (Chapter 14).

\noindent\textbf{How the book is organized.} This book can be used for a number of 
different courses. I have attempted to organize it as flexibly as possible, so as to 
enable the instructor to follow his or her own preferences in the matter.

Part I, consisting of the first eight chapters, is devoted to the subject commonly 
called general topology. The first four chapters deal with the body of material that, 
in my opinion, should be included in any introductory topology course worthy of 
the name. This may be considered the ``irreducible core'' of the subject, treating as 
it does set theory, topological spaces, connectedness, compactness (through compactness of 
finite products), and the countability and separation axioms (through the Urysohn 
metrization theorem). The remaining four chapters of Part I explore additional topics; 
they are essentially independent of one another, depending on only the core material 
of Chapters 1-4. The instructor may take them up in any order he or she chooses.

Part II constitutes an introduction to the subject of Algebraic Topology. 
It depends on only the core material of Chapters 1-4. This part of the book treats 
with some thoroughness the notions of fundamental group and covering space, along 
with their many and varied applications. Some of the chapters of Part II are 
independent of one another; the dependence among them is expressed in the following
diagram:

\begin{center}
\begin{tikzpicture}[
    >=Stealth,
    every node/.style={
        font=\normalfont,
        anchor=west,
    }
]

\def\rowsep{1.4cm}
\def\indent{2.4cm}
% 纵向箭头调整
\def\vbaseone{-0.24cm}
\def\vbasetwo{0.36cm}
\def\vbasethree{0.36cm}
\def\vbasefour{0.48cm}
\def\vbasefive{-0.12cm}
\def\vtop{0.48cm}
% 第9章箭头起点3个端点
\def\hoffsetAone{-2.16cm}
\def\hoffsetAtwo{-2.4cm}
\def\hoffsetAthree{-2.64cm}
% 第9章箭头结尾3个端点
\def\hoffsetBone{-5.0cm}
\def\hoffsetBtwo{-5.0cm}
\def\hoffsetBthree{-4.8cm}
% 第11章到第14章2个箭头末尾
\def\hoffsetCone{-4.0cm}
\def\hoffsetCtwo{-3.4cm}
% 第13章到第14章箭头
\def\hoffsetDone{-4.8cm}
\def\hoffsetDtwo{-3.0cm}

\node (c9)  at (0,0)                {Chapter 9 The Fundamental Group};
\node (c10) at (\indent, -\rowsep)   {Chapter 10 Separation Theorems in the Plane};
\node (c11) at (\indent, -2*\rowsep) {Chapter 11 The Seifert-van Kampen Theorem};
\node (c12) at (\indent, -3*\rowsep) {Chapter 12 Classification of Surfaces};
\node (c13) at (\indent, -4*\rowsep) {Chapter 13 Classification of Covering Spaces};
\node (c14) at (0, -5*\rowsep)       {Chapter 14 Applications to Group Theory};

\draw[->] ($(c9.base) + (\hoffsetAone, \vbaseone)$) -- ($(c10.base) + (\hoffsetBone, \vbasetwo)$);
\draw[->] ($(c9.base) + (\hoffsetAtwo, \vbaseone)$) -- ($(c11.base) + (\hoffsetBtwo, \vbasethree)$);
\draw[->] ($(c9.base) + (\hoffsetAthree, \vbaseone)$) -- ($(c13.base) + (\hoffsetBthree, \vbasefour)$);
\draw[->] ($(c11.base) + (\hoffsetBtwo+0.2cm, \vbasefive)$) -- ($(c12.base) + (\hoffsetCone+0.23cm, \vtop)$);
\draw[->] ($(c11.base) + (\hoffsetBone, \vbasefive)$) -- ($(c14.base) + (\hoffsetCtwo, \vtop)$);
\draw[->] ($(c13.base) + (\hoffsetDone, \vbasefive)$) -- ($(c14.base) + (\hoffsetDtwo, \vtop)$);

\end{tikzpicture}
\end{center}

Certain sections of the book are marked with an asterisk; these sections may be 
omitted or postponed with no loss of continuity. Certain theorems are marked similarly. 
Any dependence of later material on these asterisked sections or theorems is indicated 
at the time, and again when the results are needed. Some of the exercises also depend 
on earlier asterisked material, but in such cases the dependence is obvious.

Sets of supplementary exercises appear at the ends of several of the chapters. 
They provide an opportunity for exploration of topics that diverge somewhat from 
the main thrust of the book; an ambitious student might use one as a basis for an 
independent paper or research project. Most are fairly self-contained, but the one 
on topological groups has as a sequel a number of additional exercises on the topic 
that appear in later sections of the book.

\noindent\textbf{Possible course outlines. }Most instructors who use this text for a 
course in general topology will wish to cover Chapters 1-4, along with the Tychonoff 
theorem in Chapter 5. Many will cover additional topics as well. Possibilities include 
the following: the Stone-Čech compactification (\S38), metrization theorems 
(Chapter 6), the Peano curve (\S44), Ascoli's theorem (\S45 and/or \S47), and 
dimension theory (\S50). I have, in different semesters, followed each of these options.

For a one-semester course in algebraic topology, one can expect to cover most of Part II.

It is also possible to treat both aspects of topology in a single semester, although 
with some corresponding loss of depth. One feasible outline for such a course would 
consist of Chapters 1-3, followed by Chapter 9; the latter does not depend on the 
material of Chapter 4. (The non-asterisked sections of Chapters 10 and 13 also are 
independent of Chapter 4.)

\noindent\textbf{Comments on this edition.} The reader who is familiar with the first 
edition of this book will find no substantial changes in the part of the book dealing 
with general topology. I have confined myself largely to ``fine-tuning'' the text 
material and the exercises. However, the final chapter of the first edition, which 
dealt with algebraic topology, has been substantially expanded and rewritten. It has 
become Part II of this book. In the years since the first edition appeared, it has 
become increasingly common to offer topology as a two-term course, the first devoted 
to general topology and the second to algebraic topology. By expanding the treatment 
of the latter subject, I have intended to make this revision serve the needs of such 
a course.

\noindent\textbf{Acknowledgments. }Most of the topologists with whom I have studied, 
or whose books I have read, have contributed in one way or another to this book; 
I mention only Edwin Moise, Raymond Wilder, Gail Young, and Raoul Bott, but there 
are many others. For their helpful comments concerning this book, my thanks to Ken 
Brown, Russ McMillan, Robert Mosher, and John Hemperly, and to my colleagues George 
Whitehead and Kenneth Hoffman.

The treatment of algebraic topology has been substantially influenced by the 
excellent book by William Massey \cite{Massey1990}, to whom I express appreciation. Finally, 
thanks are due Adam Lewenberg of MacroTeX for his extraordinary skill and patience 
in setting text and juggling figures.

But most of all, to my students go my most heartfelt thanks. From them I 
learned at least as much as they did from me; without them this book would be very 
different.

% 作者信息靠右对齐
\begin{flushright}
J.R.M.
\end{flushright}